\chapter{Ergebnisse und Interpretation}
\label{ch:ergebnisse}

% Die Ergebnisse werden hier unmittelbar eingeordnet, statt die Deutung in ein
% eigenes Diskussionskapitel auszulagern. Der Leser hat die zugehoerigen
% Tabellen und Abbildungen damit an derselben Stelle vor sich.

Grundlage dieses Kapitels ist der finale Lauf: 120 Trainingsläufe über 24
Matchups und fünf Seeds, gefolgt von 9000 Evaluationsepisoden über 360 Läufe.
In allen folgenden Tabellen steht die Zeile für die \emph{verteidigte}
Topologie und die Spalte für den Angreifer, also für die Topologie, auf der
dieser trainiert wurde.

Zahlen werden in diesem Kapitel nach drei festen Vorschriften gebildet, die
hier einmal benannt seien, statt sie bei jeder Größe zu wiederholen.

\emph{Verlaufswerte} beschreiben die Form der Lernkurven in den Abbildungen:
je Episode das Mittel über die fünf Seeds, darüber ein zentriertes gleitendes
Mittel über fünf Episoden. Über die Seeds wird bewusst gemittelt und nicht
der Median genommen: Ein Einbruch, der nur einzelne Läufe betrifft, soll in
der Kurve sichtbar bleiben, und gerade solche Einbrüche werden in
\Cref{sec:lernkurven} besprochen. Die Glättung dient allein der Lesbarkeit.

\emph{Niveauwerte} beziffern, wo ein Training beginnt und wo es endet: je
Seed der Median seiner ersten beziehungsweise letzten zehn Episoden, darüber
das Mittel der fünf Seeds; sie heißen im Folgenden \emph{Startniveau} und
\emph{Endniveau}. Der Median je Seed wirkt vor allem am Trainingsende: Dort kann eine
einzelne Episode mit \ac{SLA}-Brüchen um Größenordnungen unter dem typischen
Wert liegen und würde ein Mittel allein bestimmen; am Anfang, wo alle zehn
Episoden gleichermaßen extrem ausfallen, gibt er schlicht die Mitte des
Fensters wieder. Das anschließende Mittel hält alle fünf Seeds im Ergebnis, sodass ein abweichender Lauf den
Zellwert sichtbar verschiebt. Alle Anfangs- und Endwerte in Text und
Tabellen sind solche Niveauwerte und keine abgelesenen Kurvenpunkte. Der
Grund liegt am Rand der Kurven: Das zentrierte Glättungsfenster ragt dort
über die vorhandenen Episoden hinaus, bei Episode~1 stehen von fünf
Episoden nur drei, und der abgelesene Wert hinge damit mehr an der
gewählten Fensterbreite als an den Daten. Wo der Text einen Wert aus dem
Inneren einer Kurve zitiert, ist er als Wert der Mittelkurve gekennzeichnet.

\emph{Quoten} wie die Crown-Jewel-Quote beruhen je Episode und Seed auf
einem einzelnen Ja oder Nein. Fünf Seeds erlauben je Episode damit nur sechs
mögliche Quoten in Stufen von zwanzig Prozentpunkten; ein Verlauf ist so
nicht ablesbar. Quoten werden deshalb gemeinsam über ein Fenster
von elf Episoden und alle fünf Seeds gebildet, also aus 55 Beobachtungen je
Punkt. In der Evaluation entfällt das Fenster: Die Policy ist dort fertig trainiert, es gibt keinen Verlauf mehr, und eine Quote fasst schlicht alle
125 Episoden einer Zelle zusammen.

\section{Lernkurven je Topologie}
\label{sec:lernkurven}

Die erste Frage an einen Trainingslauf ist, ob überhaupt etwas gelernt wurde.
Sie lässt sich am Verlauf des Episoden-Rewards ablesen, aber nicht an dem des
Verteidigers allein.

Dessen Kurve steigt nämlich in allen 120 Läufen, und kein einziger
verschlechtert sich. Das ist weniger ein Befund als eine Folge des Aufbaus.
Ein untrainierter Verteidiger sperrt wahllos, unterschreitet damit die
Verfügbarkeitsschwelle und kassiert die Strafe von $-50$ über hunderte
Schritte einer Episode (\Cref{sec:sla}). Die Startwerte fallen dadurch extrem
aus und sagen kaum etwas über den Angreifer; beispielsweise liegt das
Startniveau von \emph{flat} gegen den eigenen Angreifer bei $-8381$ (\Cref{tab:reward-ende}). Dass der Agent diese Strafe binnen weniger Episoden abstellt,
ist die erste und drastischste Korrektur; eine Verbesserung ist von
diesem Startniveau aus fast unvermeidlich.

An der zweiten Kurve entscheidet sich deshalb die eigentliche Frage: Gelingt es dem
Verteidiger, den Reward des Angreifers zu senken, ihn also tatsächlich am
Vordringen zu hindern? Nur diese Größe entsteht allein aus dem Geschehen im
Netz und ist von Abwehrbonus und \ac{SLA}-Strafe unberührt
(\Cref{sec:rewards}).

Auch darauf lautet die Antwort ja, allerdings in sehr unterschiedlichem
Ausmaß. \Cref{tab:reward-atk-verlauf} stellt Anfang und Ende für alle 24
Matchups gegenüber.

\begin{table}[H]
  \centering
  \caption[Angreifer-Reward als Start- und Endniveau]{Angreifer-Reward als
  Start- und Endniveau. Hervorgehoben sind die Zellen, in denen Verteidiger
  und Angreifer auf derselben Topologie entstanden sind.}
  \label{tab:reward-atk-verlauf}
  \footnotesize
  \begin{tabular}{l*{6}{r}}
    \toprule
    Verteidigt & Chain & Flat & H\,\&\,S & DMZ & Micro & Super \\
    \midrule
    Flat        & 296\,$\to$\,160 & \diag{1349\,$\to$\,664} & 325\,$\to$\,164 & 257\,$\to$\,147 & 300\,$\to$\,166 & 581\,$\to$\,188 \\
    Hub \& Spoke& 230\,$\to$\,110 & 271\,$\to$\,62 & \diag{758\,$\to$\,75} & 139\,$\to$\,80 & 253\,$\to$\,108 & 275\,$\to$\,64 \\
    DMZ         & 208\,$\to$\,80 & 161\,$\to$\,66 & 225\,$\to$\,86 & \diag{901\,$\to$\,150} & 137\,$\to$\,64 & 164\,$\to$\,79 \\
    Micro-Seg.  & 177\,$\to$\,77 & 244\,$\to$\,75 & 194\,$\to$\,49 & 68\,$\to$\,52 & \diag{655\,$\to$\,20} & 495\,$\to$\,77 \\
    \bottomrule
  \end{tabular}
\end{table}

Die hervorgehobenen Zellen sind diejenigen, in denen die verteidigte Topologie
und die Trainingstopologie des Angreifers übereinstimmen. Da die Spalten in
derselben Reihenfolge stehen wie die Zeilen, bilden sie in der Tabelle eine
Diagonale; dieser Begriff wird im Folgenden für sie verwendet.

Die seedweise Bildung der Niveauwerte ist dabei keine Förmlichkeit. Da das Budget in
Zeitschritten und nicht in Episoden vergeben ist, umfassen die Läufe
unterschiedlich viele Episoden. Deutlich tritt das in drei Matchups auf: 251
bis 253 in \emph{flat} gegen \emph{flat}, 252 bis 259 in \emph{hub\_and\_spoke} gegen \emph{hub\_and\_spoke}
und 261 bis 311 in \emph{dmz} gegen \emph{dmz}. Würde man erst über die Seeds mitteln und
dann die letzten Episoden nehmen, so beruhte das Ende einer solchen Zelle
auf immer weniger Läufen. Alle drei betroffenen Matchups liegen auf der
Diagonalen; eine ungenauere Bildung hätte also gerade die Zellen verzerrt,
an denen die Befunde dieses Kapitels hängen.

Schon die Startwerte zeigen, dass die Herkunft eines Angreifers zählt. Gegen
einen noch untrainierten Verteidiger erreicht in jeder der vier Zeilen
derjenige Angreifer den höchsten Reward, der auf ebendieser Topologie
trainiert wurde. Der Abstand zum Zweitbesten reicht dabei vom Faktor 1,3 in
\emph{micro\_segmented} bis zum Faktor 4,0 in \emph{dmz}, wo der Spezialist 901 Punkte
erzielt und der nächstbeste Angreifer 225. Zweitbester ist in \emph{flat} und
\emph{micro\_segmented} jeweils der Super-Angreifer, also der einzige, der ebenfalls
auf dieser Topologie trainiert hat, wenn auch nicht ausschließlich.

Außerhalb der Diagonalen behält der Angreifer am Ende zwischen 16 und
76~Prozent seines anfänglichen Rewards, im Mittel 42~Prozent. Auf der
Diagonalen bricht er dagegen auf 3 bis 17~Prozent ein. Der Angreifer, der die
verteidigte Struktur kennt, startet also von einem weit höheren Niveau und
wird dennoch am gründlichsten zurückgedrängt.

Eine Zelle folgt diesem Muster nicht. \emph{flat} gegen \emph{flat} startet mit 1349 ebenfalls
hoch, endet aber bei 664 und behält
damit knapp die Hälfte, also ähnlich viel wie ein gewöhnliches Matchup. Dieser Fall wird das gesamte Kapitel begleiten.

Für jede verteidigte Topologie zeigt eine eigene Abbildung den Verlauf
beider Agenten gegen alle sechs Angreifer. Jedes der sechs Teilbilder
trägt dabei eine eigene Skala, denn die Anfangswerte der einzelnen Matchups liegen um mehr
als eine Größenordnung auseinander.

\begin{figure}[H]
  \centering
  \includegraphics[width=\textwidth]{lernkurven_flat}
  \caption[Lernkurven auf \emph{flat}]{Verteidiger auf \emph{flat} gegen alle
  sechs Angreifer, gemittelt über fünf Seeds und geglättet über das gleitende
  Mittel von fünf Episoden.}
  \label{fig:lk-flat}
\end{figure}

Die Zeile \emph{flat} in \Cref{tab:reward-atk-verlauf} zeigt den Angreifer in fünf
der sechs Spalten auf einem höheren Endniveau als in jeder anderen
verteidigten Topologie; einzig der DMZ-Angreifer endet auf seiner
Heimattopologie mit 150 knapp über den 147. In \Cref{fig:lk-flat} kommt
entsprechend keine der sechs Verteidigerkurven der Nulllinie nennenswert
nahe. Am deutlichsten ist das im Teilbild des eigenen Angreifers, oben in der
Mitte; \Cref{tab:reward-ende} beziffert es: Startniveau $-8381$, Endniveau
$-710$, während dieselbe Topologie gegen alle übrigen Angreifer zwischen
$-168$ und $-197$ endet.

\begin{figure}[H]
  \centering
  \includegraphics[width=\textwidth]{lernkurven_hub_and_spoke}
  \caption[Lernkurven auf \emph{hub\_and\_spoke}]{Verteidiger auf
  \emph{hub\_and\_spoke} gegen alle sechs Angreifer.}
  \label{fig:lk-hs}
\end{figure}

In \emph{hub\_and\_spoke} (\Cref{fig:lk-hs}) fällt das Teilbild des eigenen
Angreifers durch etwas anderes auf. Dort steigt die Mittelkurve des Angreifers
zunächst bis auf über $+3300$ in Episode~25, also weit über das Niveau aller
übrigen Teilbilder, und fällt dann über das Training fast vollständig wieder ab. Der Verteidiger beginnt im Startniveau bei
$-7910$ und endet bei $-78$ (\Cref{tab:reward-ende}). Der Anfang ist dabei kein
Spiegelbild: Der Verteidiger verliert ein Vielfaches dessen, was der Angreifer
gewinnt, weil zur Kopplung noch die \ac{SLA}-Strafe hinzukommt. Erst am Ende
stehen die beiden Werte einander spiegelbildlich gegenüber, und genau das
zeigt an, dass keine Brüche mehr anfallen. Der auf dieser Topologie trainierte
Angreifer beginnt also mit Abstand am stärksten und verliert im Verlauf
90~Prozent seines Rewards. Am Ende ist er von den übrigen fünf Angreifern
dieser Zeile nicht mehr zu unterscheiden, die zwischen 62 und 110 liegen
(\Cref{tab:reward-atk-verlauf}).

In dieser Abbildung fallen zudem die tiefen Einbrüche auf, etwa gegen
Episode~150 im selben Teilbild und gegen Episode~240 gegen den
Micro-Segmentation-Angreifer. Beide gehen auf einzelne Seeds zurück und haben
einen unterschiedlichen Umfang. Der erste ist eine Phase: In zwei der fünf
Läufe häufen sich zwischen Episode~110 und 165 die \ac{SLA}-Brüche, in einem
über 21 und im anderen über 16 Episoden. Der zweite ist eine einzige Episode,
in der die Verfügbarkeit über \num{1668} der rund 2000 Schritte unter der
Schwelle lag.

\begin{figure}[H]
  \centering
  \includegraphics[width=\textwidth]{sla_einbrueche}
  \caption[Reward und \acs{SLA}-Brüche je Seed]{Verteidiger-Reward und Schritte
  mit \ac{SLA}-Bruch für die beiden auffälligen Matchups, je Seed einzeln
  aufgetragen und nicht gemittelt. Der Reward ist über fünf Episoden geglättet,
  die Bruchschritte sind es nicht.}
  \label{fig:sla-einbrueche}
\end{figure}

\Cref{fig:sla-einbrueche} stellt beides nebeneinander. Die Ausschläge der
Reward-Kurve und die Spitzen der Bruchschritte liegen an denselben Episoden,
und sie beschränken sich auf wenige Läufe: links auf zwei, von denen einer
deutlich stärker ausschlägt, rechts auf einen einzigen. Kurze Brüche
treten auch in den übrigen Seeds vereinzelt auf; Episoden mit mehreren hundert
Bruchschritten bleiben jedoch auf diese wenigen Läufe beschränkt. Im gesamten
Training gibt es nur zwölf Episoden mit mehr als 500 Bruchschritten, und acht
davon liegen in den ersten beiden Episoden ihres Laufs.

Eine Erklärung liefert der Vergleich innerhalb des Laufs mit der schwersten
Episode, also des einen betroffenen Seeds gegen den
Micro-Segmentation-Angreifer; als ruhig gelten dabei seine bruchfreien
Episoden ab Episode~51, nach der Anfangsphase.
In der Episode mit \num{1668} Bruchschritten bleibt die Zahl
der Reimages mit 117 gegenüber 112 im Median praktisch unverändert. Die Zahl
der Sperren steigt dagegen von 7 auf 36, und keine einzige von ihnen wird
wieder aufgehoben, gegenüber im Median 61 Freigaben in einer ruhigen Episode. Nicht das
Neuaufsetzen bricht also die Verfügbarkeit, sondern Sperren, die über den
größten Teil der Episode Bestand haben; eine gesetzte Sperre wirkt bis zu
ihrer Aufhebung fort. Eine erreichte
Policy ist demnach nicht dauerhaft stabil, und das Konvergenzkriterium aus
\Cref{sec:konvergenz} erfasst solche späten Rückfälle nicht.

Woher der Umschwung in der Policy rührt, lässt sich aus den erhobenen Daten
nicht eindeutig bestimmen. Zweierlei ist dazu festzuhalten. Erstens ist es kein
allgemeines Verhalten des Verfahrens, sondern eines einzelner Läufe: Von fünf
Seeds sind jeweils einer oder zwei betroffen, die übrigen zeigen keine solchen
Ausschläge in der Zahl der Bruchschritte.
Genau diese Streuung zwischen sonst gleichen Läufen beschreiben Henderson
et~al.\ als eigenständige Schwierigkeit von Policy-Gradient-Verfahren; sie
zeigen, dass allein die Wahl des Zufallsstartwerts zu deutlich verschiedenen
Ergebnisverteilungen führen kann und ein Mittel über wenige Läufe deshalb
keine belastbaren Schlüsse
erlaubt~\cite{henderson2019deepreinforcementlearningmatters}. Für die vorliegende Arbeit begrenzt das lediglich die Deutung: Die
Einbrüche in den Mittelkurven sind keine Aussage über die Topologie,
sondern die sichtbare Spur einzelner Seeds.

Zweitens ist der naheliegende Verdacht, eine einzelne Ausreißerepisode habe
die Policy unmittelbar verschoben, mit dem Verfahren nicht vereinbar. Genau
das verhindert die Begrenzung aus \Cref{eq:ppo-clip}: Sie deckelt, wie weit
sich die Policy je Aktualisierung von der bisherigen entfernen darf,
unabhängig davon, wie groß der Vorteil ausfällt.

Diese Grenze gilt allerdings nur für die Policy selbst. Der Critic darf
seine Prognose nach einer solchen Episode beliebig stark anpassen, und da
sich beide in der eingesetzten Umsetzung die vorderen Schichten des Netzes
teilen, wirkt seine Korrektur mittelbar auch auf den Actor -- an dessen
eigener Begrenzung vorbei. Derselbe Mechanismus trägt die rasche Rückkehr:
Sobald die Episoden wieder normal ausfallen, fallen die Schätzungen des
Critic deutlich negativer aus als das tatsächliche Geschehen. Jede
gewöhnliche Episode erscheint dadurch besser als erwartet, das zuletzt
gezeigte Verhalten wird verstärkt, und die Policy findet auf das alte
Niveau zurück. Was den Ausschlag überhaupt auslöst, erklären diese Überlegungen
nicht.

\begin{figure}[H]
  \centering
  \includegraphics[width=\textwidth]{lernkurven_dmz}
  \caption[Lernkurven auf \emph{dmz}]{Verteidiger auf \emph{dmz} gegen alle
  sechs Angreifer.}
  \label{fig:lk-dmz}
\end{figure}

\emph{dmz} (\Cref{fig:lk-dmz}) zeigt dasselbe Muster, aber deutlich langsamer.
Die Mittelkurve des Angreifers steigt zunächst noch, bis auf 4257 in
Episode~22, und geht erst danach zurück: auf 2277 in Episode~75, 946 in
Episode~100 und 282 in Episode~110; von dort an ändert sich wenig mehr. Die
Mittelkurve des Verteidigers verbessert sich über denselben Zeitraum stetig
und liegt in Episode~50 bereits bei $-1030$. Da der Angreifer
eingefroren ist, lernt er selbst nichts dazu; eine gerichtete Veränderung
seiner Kurve geht deshalb auf den Verteidiger zurück, während kurzfristige
Schwankungen der Streuung zwischen den Läufen zuzuschreiben sind.
Erst nach rund 25 Episoden beginnt dieser, ihn tatsächlich zu
bremsen, und dafür braucht er dann etwa 80 weitere. In demselben
Zeitraum bricht auch die Crown-Jewel-Quote dieses Matchups ein
(\Cref{fig:cj-verlauf-dmz}); beide Größen bewegen sich also gemeinsam.

\begin{figure}[H]
  \centering
  \includegraphics[width=\textwidth]{cj_verlauf_dmz}
  \caption[Reward und Crown-Jewel-Quote in \emph{dmz}]{Angreifer-Reward und
  Crown-Jewel-Quote im Matchup \emph{dmz} gegen \emph{dmz}.}
  \label{fig:cj-verlauf-dmz}
\end{figure}

Bis Episode~55 fällt das Crown Jewel in jeder einzelnen Episode. Danach
sinkt die Quote parallel zum Reward, auf 85~Prozent in Episode~100 und auf
ihren Tiefstwert von 26~Prozent in Episode~145. Der Rückgang des Rewards
ist also kein bloßer Punktverlust, sondern bedeutet, dass der Verteidiger
den Angreifer in einem wachsenden Teil der Episoden vom wertvollsten
System fernhält.

Danach pendelt die Quote zwischen 26 und 62~Prozent, ohne weiter zu
fallen; die noch höheren Werte am rechten Rand beruhen bereits auf weniger
als fünf Läufen. Das ist jedoch keine gemeinsame Bewegung aller Läufe,
sondern deren Streuung: Über die Episoden~150 bis 250 erreichen die fünf
Seeds Quoten von 26, 31, 37, 59 und 70~Prozent, jeder für sich weitgehend
stabil. Die gemittelte Kurve bildet damit eine Mischung aus fünf
verschiedenen, aber jeweils gefestigten Policies ab, nicht einen
Rückfall. Die Evaluation bestätigt das Niveau später mit 36~Prozent
(\Cref{tab:cj-eval}).

\begin{figure}[H]
  \centering
  \includegraphics[width=\textwidth]{lernkurven_micro_segmented}
  \caption[Lernkurven auf \emph{micro\_segmented}]{Verteidiger auf
  \emph{micro\_segmented} gegen alle sechs Angreifer.}
  \label{fig:lk-ms}
\end{figure}

In \emph{micro\_segmented} (\Cref{fig:lk-ms}) verlaufen beide Kurven im Teilbild
des eigenen Angreifers über das gesamte Training aufeinander zu und enden bei den Niveaus $+20$
und $-20$. Anders als in \emph{dmz} beginnt der Rückgang
sofort: Die Mittelkurve fällt auf 363 in Episode~25 und 193 in Episode~50,
danach langsamer weiter. Von den vier Matchups auf der Diagonalen ist dies das einzige, in dem der
Angreifer am Ende praktisch nichts mehr erwirtschaftet. Innerhalb der
Abbildung ähnelt dem Teilbild des eigenen Angreifers einzig das des
Super-Angreifers: Dessen Mittelkurve liegt in Episode~25 mit
208 noch weit über den vier übrigen Teilbildern, die dort 58 bis 127 zeigen,
und die Verteidigerkurve beginnt gegen ihn tiefer und holt langsamer auf --
in Episode~25 steht sie bei $-494$, in den vier übrigen Teilbildern zwischen
$-72$ und $-228$. Das fügt sich in das Bild, denn der Super-Angreifer ist
der einzige weitere, der diese Topologie im Training gesehen hat, als letzte Station seines gestuften Trainings;
\Cref{sec:angreiferherkunft} nimmt das wieder auf.
Die Anfangswerte fallen zugleich am mildesten aus: Das tiefste Startniveau
liegt hier bei $-3711$ gegenüber $-8381$ in \emph{flat}
(\Cref{tab:reward-ende}). Das passt zu den nur 145 erreichbaren Punkten dieser Topologie.

Über alle vier Abbildungen hinweg zeigt sich damit dasselbe Muster: Das Teilbild
des eigenen Angreifers ist stets dasjenige mit dem tiefsten Startwert und dem
längsten Weg. Ein Verteidiger tut sich gegen den Gegner am schwersten, der
seine eigene Struktur kennt. \Cref{sec:angreiferherkunft} geht dem nach.

\Cref{tab:reward-ende} stellt Start- und Endniveau des Verteidigers
gegenüber. Am Start trägt die Diagonale in jeder Zeile mit Abstand den
tiefsten Wert; bemerkenswert ist, dass ihre Werte am Ende nicht mehr
herausstechen:
In \emph{micro\_segmented} ist sie sogar der beste Wert der Zeile, in
\emph{hub\_and\_spoke} liegt sie im Mittelfeld. Der anfängliche Rückstand gegen den
passgenauen Angreifer ist bis zum Trainingsende also aufgeholt. Die Ausnahme
ist erneut \emph{flat} gegen \emph{flat} mit $-710$, dem mit Abstand schlechtesten Wert der
Tabelle; alle übrigen 23 Zellen liegen zwischen $-20$ und $-224$. Der
zweitschlechteste Wert, $-224$ gegen den Hub-and-Spoke-Angreifer in \emph{dmz}, wird
dabei von einem einzelnen Seed bei $-839$ getragen; die übrigen vier Läufe
dieser Zelle enden zwischen $-58$ und $-86$.

\begin{table}[H]
  \centering
  \caption[Verteidiger-Reward als Start- und Endniveau]{Start- und Endniveau
  des Verteidiger-Rewards. Werte näher an null sind besser.}
  \label{tab:reward-ende}
  \footnotesize
  \setlength{\tabcolsep}{4pt}
  \begin{adjustbox}{max width=\textwidth}
  \begin{tabular}{l*{6}{r}}
    \toprule
    Verteidigt & Chain & Flat & H\,\&\,S & DMZ & Micro & Super \\
    \midrule
    Flat        & $-1118\to-182$ & \diag{$-8381\to-710$} & $-1694\to-193$ & $-1209\to-168$ & $-1318\to-188$ & $-2349\to-197$ \\
    Hub \& Spoke& $-983\to-109$ & $-951\to-64$ & \diag{$-7910\to-78$} & $-789\to-81$ & $-946\to-117$ & $-1600\to-72$ \\
    DMZ         & $-1159\to-77$ & $-977\to-67$ & $-1193\to-224$ & \diag{$-7265\to-157$} & $-978\to-62$ & $-866\to-76$ \\
    Micro-Seg.  & $-735\to-76$ & $-1156\to-74$ & $-904\to-49$ & $-434\to-53$ & \diag{$-3711\to-20$} & $-2799\to-83$ \\
    \bottomrule
  \end{tabular}
  \end{adjustbox}
\end{table}

Diese Zahlen sind allerdings nur innerhalb einer Zeile vergleichbar, nicht
zwischen den Zeilen. Der Verteidiger-Reward ist über die Grundkopplung an den
des Angreifers gebunden, und der erreichbare Punktevorrat unterscheidet sich
je Topologie: In \emph{micro\_segmented} sind vom Einstiegsknoten aus nur 145
der 235 Punkte überhaupt erreichbar (\Cref{sec:topomodell}). Ein Teil des
besseren Rewards geht dort also allein darauf zurück, dass weniger zu holen
war; welcher Anteil das ist, lässt sich der Zahl nicht entnehmen. Die Aussagen zur
Wirksamkeit stützen sich deshalb in \Cref{sec:wirksamkeitergebnis} auf
Größen, die von der Reward-Skala unabhängig sind.

\section{Crown-Jewel-Quote im Trainingsverlauf}
\label{sec:cj-verlauf}

Die Reward-Kurven laufen in allen 24 Matchups früh auf ein Niveau zu, auf dem
sie anschließend verharren. Daraus ließe sich schließen, dass auch die
Verteidigung ab diesem Punkt nicht mehr besser wird. Die Crown-Jewel-Quote
widerspricht dem. Sie ist zugleich die Größe, an der die Fragestellung
unmittelbar hängt, denn sie sagt, ob der Angreifer das wertvollste System
erreicht.

\begin{figure}[H]
  \centering
  \includegraphics[width=\textwidth]{cj_verlauf_diagonale}
  \caption[Crown-Jewel-Quote im Trainingsverlauf]{Anteil der Episoden mit
  erreichtem Crown Jewel für die vier Matchups auf der Diagonalen. Die Quote
  ist über ein Fenster von elf Episoden und alle fünf Seeds gebildet, also über
  55~Episoden je Punkt.}
  \label{fig:cj-diagonale}
\end{figure}

\Cref{fig:cj-diagonale} trennt die vier Topologien deutlicher als jede andere
Größe dieses Kapitels. Zu Beginn erreicht der Angreifer das Ziel überall in
jeder Episode. Danach gehen die Verläufe auseinander:

\begin{itemize}
  \item In \emph{flat} verhindert der Verteidiger über das gesamte Training
    keinen einzigen Zugriff auf das Crown Jewel, von einer einzigen Episode
    abgesehen.
  \item In \emph{micro\_segmented} fällt sie stetig und liegt kurz nach
    Episode~230 bei null; über die letzten zwanzig Episoden aller fünf
    Seeds erreicht der Angreifer das Ziel in keiner einzigen.
  \item \emph{hub\_and\_spoke} und \emph{dmz} liegen dazwischen und kommen auf
    28 beziehungsweise 42~Prozent in den letzten zwanzig Episoden.
\end{itemize}

Bemerkenswert ist der Zeitpunkt. In \emph{micro\_segmented} sinkt die Quote
zwischen Episode~150 und 230 von 38 auf 5~Prozent, während der Reward in
diesem Abschnitt nur noch von etwa $-100$ auf $-20$ steigt. Am Reward gemessen
geschieht dort fast nichts mehr, an der Crown-Jewel-Quote gemessen das
Entscheidende. 
Der Grund liegt in der Reward-Konstruktion: Ein Zugriff, der
verhindert wird, schlägt sich nur mit dem Wert des Zielknotens von
100~Punkten nieder, weil dem Angreifer nichts genommen, sondern nur nichts
gegeben wird (\Cref{sec:marlon-original}).

In \emph{dmz} erreicht der Angreifer das Crown Jewel auch nach dem
Einpendeln regelmäßig, ohne dass sich das im Reward niederschlägt
(\Cref{fig:cj-verlauf-dmz}, bei den Lernkurven). Der Grund liegt in der
Kopplung aus Verlust und Rückeroberung (\Cref{sec:rewards}). Über die
Episoden~150 bis 250 verliert der Angreifer in den Episoden mit
erreichtem Crown Jewel im Median 2035~Punkte durch neu aufgesetzte Knoten,
gegenüber 1552 in den übrigen. Er nimmt das Ziel also,
verliert es wieder und nimmt es erneut, und jede dieser Runden gleicht sich
aus. Am Ende trennt die Episoden mit erreichtem Crown
Jewel von denen ohne ein Unterschied von 175 zu 126~Punkten -- weniger als
die 100~Punkte, die das gehaltene Crown Jewel allein wert wäre.

Beide Größen messen demnach Verschiedenes: Die Crown-Jewel-Quote sagt, ob der
Angreifer das Ziel überhaupt erreicht, der Reward, ob er es halten kann.
Sie bewegen sich dennoch in dieselbe Richtung. Ordnet man die fünf Seeds nach ihrer Quote, so
steigt der Reward-Median mit ihr, von 123 bei 26~Prozent bis 188 bei
70~Prozent. Erst aus beiden zusammen ergibt sich ein Bild, wie in
\Cref{sec:metriken} angelegt.

\section{Gehaltene Knoten im Trainingsverlauf}
\label{sec:knoten-verlauf}

Ob das Crown Jewel gefallen ist, beantwortet noch nicht, wie weit der Angreifer
insgesamt vorgedrungen ist. Ein Netz, in dem er acht Knoten hält, und eines, in
dem er auf dem Einstiegsknoten festsitzt, stehen in dieser Quote gleich, sofern
das Ziel in beiden Fällen unangetastet blieb. Dafür dient die dritte Zielgröße
aus \Cref{sec:metriken}, die größte Zahl gleichzeitig kompromittierter Knoten
je Episode. Der Einstiegsknoten zählt dabei mit: Ein Wert von 1,0 bedeutet, dass
der Angreifer über ihn nicht hinausgekommen ist.

\begin{figure}[H]
  \centering
  \includegraphics[width=\textwidth]{knoten_verlauf_diagonale}
  \caption[Gehaltene Knoten im Trainingsverlauf]{Größte Zahl gleichzeitig
  gehaltener Knoten je Episode für die vier Matchups auf der Diagonalen,
  gemittelt über fünf Seeds und geglättet über fünf Episoden. Von zehn Knoten
  im Netz.}
  \label{fig:knoten-diagonale}
\end{figure}

In \emph{dmz} und \emph{hub\_and\_spoke} steigt die Mittelkurve zunächst noch
an und erreicht 9,6 beziehungsweise 9,4. Der Angreifer hält dort in den meisten
Episoden das gesamte Netz. Erst danach setzt der Rückgang ein.
\emph{micro\_segmented} beginnt dagegen bei 4,0 und hält damit ebenfalls
alles, was von dort aus überhaupt erreichbar ist (\Cref{sec:topomodell}).

Die Abbildung zeigt nur die Diagonale, alle 24 Matchups zeigt
\Cref{app:knoten}. Zieht man sie alle heran, so
liegen 21 von ihnen am Ende des Trainings zwischen 1,8 und 2,5. Der
Verteidiger drückt den Angreifer also im Regelfall auf ein bis zwei eroberte
Rechner zurück. Drei Zellen weichen davon ab, und zwei davon sind in der
Abbildung zu sehen: \emph{dmz} gegen \emph{dmz} mit 3,4 und \emph{flat}
gegen \emph{flat} mit 7,5, dazu \emph{flat} gegen den Super-Angreifer mit
3,0.

Auffällig ist davon nur \emph{flat} gegen \emph{flat}: Auch dort hält der
Angreifer anfangs fast das gesamte Netz, doch anders als in \emph{dmz} und
\emph{hub\_and\_spoke} setzt der Rückgang nie ein -- die Kurve bleibt über
das gesamte Training nahe ihrem Höchstwert, der Angreifer behält dauerhaft
rund drei Viertel des Netzes.

\section{Konvergenz}
\label{sec:konvergenzergebnis}

Das in \Cref{sec:konvergenz} beschriebene Kriterium wird von allen 120 Läufen
erfüllt, im Median in Episode~42. Der früheste überhaupt mögliche Wert ist
Episode~39, der früheste gemessene Episode~40. Die Mediane je verteidigter Topologie
liegen zwischen 41 und 42,5 und unterscheiden sich damit nicht sonderlich
voneinander.
Sämtliche Läufe gelten also nahezu unmittelbar an der unteren Grenze des
Kriteriums als konvergiert, obwohl sie bis zur Episode~250 laufen.

Dieser Befund ist kein Ergebnis über Topologien, sondern eines über das
Kriterium. Es misst die Restunruhe der letzten Episoden im Verhältnis zur
bisherigen Reward-Spanne des Laufs. Genau diese Spanne wird durch die ersten Episoden sehr groß, die je nach
Lauf im fünfstelligen Bereich liegen, und jede spätere
Verbesserung wirkt daran gemessen klein. Ein Lauf gilt deshalb bereits als
ruhig, während er sich tatsächlich noch verbessert.

\emph{dmz} gegen \emph{dmz} zeigt das exemplarisch: Die Spanne wird von der ersten
Episode bestimmt, die je nach Seed zwischen rund $-15\,000$ und
$-31\,000$ liegt; fünf Prozent davon sind 750 bis 1550 Punkte. Zu dem Zeitpunkt, an dem das Kriterium zutrifft, im Median in
Episode~49, steht die Mittelkurve erst
bei rund $-1000$; bis zum Trainingsende verbessert sich der Reward weiter,
auf Endniveaus zwischen $-133$ und $-196$ je Seed. Diese Verbesserung
verteilt sich über zweihundert Episoden und bleibt damit je
Fenster unter der Schwelle. Das Kriterium trennt das Ende der Unruhe nicht vom Ende
der Verbesserung.

Eine Spur des in \Cref{sec:lernkurven} beschriebenen Musters bleibt dennoch
erkennbar. In drei der vier Zeilen trägt die Diagonale den höchsten Wert, in
\emph{micro\_segmented} teilt sie ihn mit der Zelle gegen den
Super-Angreifer: \emph{dmz} gegen \emph{dmz} konvergiert im Median erst in
Episode~49, \emph{hub\_and\_spoke} gegen sich selbst in 47, \emph{micro\_segmented} in 46 und
\emph{flat} in 44, gegenüber 40
bis 46 in allen übrigen Zellen (\Cref{tab:konvergenz-vergleich}, jeweils
der Wert vor dem Schrägstrich). Die Richtung stimmt
also mit den übrigen Größen überein, die Spreizung ist aber so gering, dass
sich darauf keine Aussage stützen lässt.

Dass die Topologien sich nicht trennen, könnte an der Wahl der Schwelle liegen.
Das Kriterium besteht aus zwei voneinander unabhängigen Bedingungen, die
Verschiedenes messen: Die Streuung erfasst das Rauschen der Kurve, der Trend
erfasst, ob die Verbesserung aufgehört hat. Beide erhalten in
\Cref{sec:konvergenz} denselben Wert von 0,05 der Spanne, was eine Konvention
ist und keine Notwendigkeit.

Nachgerechnet zeigt sich, dass davon nur eine der beiden Bedingungen wirksam
ist. Die Streuung ist in fast allen 120 Läufen von dem Moment an erfüllt, ab
dem das Kriterium überhaupt geprüft wird; scheitert eine Prüfung vor der
Konvergenz, dann in aller Regel am Trend. Was den Zeitpunkt bestimmt, ist also
allein der Trendtest, und dessen Toleranz von 5~Prozent der Gesamtspanne ist weit
gefasst.

Als Gegenprobe wurde die Schwelle des Trends deshalb auf 0,01 gesenkt, die der
Streuung blieb bei 0,05. Am Kriterium selbst ändert das nichts, nur an einer
seiner beiden Zahlen.

\begin{table}[H]
  \centering
  \caption[Konvergenzepisoden je Matchup]{Konvergenzepisode je Matchup als
  Median der fünf einzeln geprüften Seeds; vor dem Schrägstrich beim
  Trendtest mit Schwelle 0,05, dahinter mit 0,01. Hervorgehoben ist die
  Diagonale.}
  \label{tab:konvergenz-vergleich}
  \footnotesize
  \setlength{\tabcolsep}{4pt}
  \begin{tabular}{l*{6}{r}}
    \toprule
    Verteidigt & Chain & Flat & H\,\&\,S & DMZ & Micro & Super \\
    \midrule
    Flat        & 41\,/\,58 & \diag{44\,/\,72} & 43\,/\,55 & 41\,/\,59 & 43\,/\,58 & 42\,/\,76 \\
    Hub \& Spoke& 41\,/\,50 & 41\,/\,55 & \diag{47\,/\,88} & 40\,/\,51 & 41\,/\,50 & 44\,/\,53 \\
    DMZ         & 41\,/\,49 & 42\,/\,53 & 41\,/\,59 & \diag{49\,/\,95} & 41\,/\,59 & 42\,/\,68 \\
    Micro-Seg.  & 40\,/\,57 & 42\,/\,49 & 41\,/\,50 & 40\,/\,52 & \diag{46\,/\,66} & 46\,/\,60 \\
    \bottomrule
  \end{tabular}
\end{table}

Auch mit dem strengeren Trendtest erfüllen alle 120 Läufe das Kriterium; der
Median wandert von Episode~42 auf 58,5. Die Topologien trennen sich weiterhin
nicht: Ihre Mediane liegen zwischen 53 und 62, die Quartile der 30 Läufe je
Topologie reichen von 55 bis 77 in \emph{flat} und von 50 bis 62 in
\emph{micro\_segmented} und überlappen damit fast vollständig.

Sichtbar wird dagegen etwas anderes. Betrachtet man nicht die Zeilen
gegeneinander, sondern die Zellen innerhalb einer Zeile, so tritt die
Diagonale hervor (\Cref{tab:konvergenz-vergleich}, Werte nach dem
Schrägstrich). Ihre Mediane liegen zwischen Episode~66 und 95, die der übrigen
20~Zellen zwischen 49 und 76. In drei der vier Zeilen ist die Diagonale die Zelle mit der spätesten
Konvergenz, in \emph{flat} die zweitspäteste hinter der Zelle gegen den
Super-Angreifer. Beim Kriterium aus \Cref{sec:konvergenz} war dieser
Unterschied kaum erkennbar. Der Befund fügt sich damit in die übrigen Größen
ein, die in \Cref{sec:angreiferherkunft} zusammengeführt werden.

Angegeben ist hier der Median der fünf Seeds einer Zelle, nicht ihr Mittel.
Einzelne Läufe konvergieren sehr viel später als die übrigen vier ihrer Zelle,
etwa \emph{flat} gegen den Hub-and-Spoke-Angreifer mit 48, 49, 55, 63 und 164.
Das Mittel dieser Zelle läge bei 76 und damit über der Diagonalen, obwohl vier
ihrer fünf Läufe zu den schnellsten der Zeile gehören.

Für die Frage nach der Topologie bleibt es beim Ergebnis: Bei fünf Seeds je
Zelle ist ein Einfluss auf die Lerngeschwindigkeit nicht nachweisbar.
Ausgeschlossen ist er damit nicht, die Daten tragen den Nachweis nur nicht.

Als Maß der Lerneffizienz ist der Konvergenzzeitpunkt in dieser Form deshalb
nur eingeschränkt geeignet. Dass die Läufe ihr volles Schrittbudget erhielten
und nicht beim Erfüllen des Kriteriums abgebrochen wurden
(\Cref{sec:konvergenz}), erweist sich hier als nützlich: Beide Varianten ließen
sich nachträglich auf vollständige Kurven anwenden, und der Vergleich in
\Cref{tab:konvergenz-vergleich} wurde dadurch überhaupt erst möglich. Wäre wie
im Referenzlauf abgebrochen worden, stünde am Ende eine einzelne Zahl, deren
Zustandekommen sich nicht mehr prüfen ließe.

Der Vergleich der Topologien stützt sich in dieser Arbeit deshalb auf das
erreichte Niveau (\Cref{sec:lernkurven} bis \Cref{sec:knoten-verlauf}) und auf
die Wirksamkeit der fertig trainierten Policy
(\Cref{sec:wirksamkeitergebnis}). \Cref{sec:limitationen} schlägt ein Maß vor,
das ohne die Normierung auf die Gesamtspanne auskommt.

\section{Aktionsverhalten des Verteidigers}
\label{sec:aktionen}

Über alle 120 Läufe gemittelt sinkt der Anteil der Sperren von 16~Prozent in
den ersten 25 Episoden auf 2~Prozent ab Episode~101. Der Verteidiger gibt das
Sperren also überwiegend auf, jedoch nicht vollständig: In
87,8~Prozent aller Episoden ist zu jedem Zeitpunkt
mindestens ein Knoten aktiv abgeschirmt. Zurück geht die Häufigkeit, mit der
Sperren umgesetzt werden, nicht ihr Bestand.

Zugleich steigt der Anteil ungültiger Aktionen von 38 auf 62~Prozent. Eine
ungültige Aktion kostet nichts und überspringt den Zug
(\Cref{sec:marlon-original}); sie ist damit der billigste Weg, nichts zu tun.
Der Aktionsraum des Verteidigers enthält keine ausdrückliche Aktion dafür.
Dasselbe gilt für Freigaben auf Knoten, die ohnehin nicht gesperrt sind: Sie
sind folgenlos. Freigaben machen über alle Läufe hinweg 36,8~Prozent aller Aktionen aus
und stehen damit nach den ungültigen Aktionen an zweiter Stelle.
Beide Anteile sind deshalb nicht als Verhalten zu lesen,
sondern zum größten Teil als Untätigkeit; einzelne Freigaben mögen berechtigt sein, aus den erhobenen Zählungen
lässt sich das jedoch nicht trennen.
Aussagekräftig bleiben Sperren und Reimages sowie ihr Verhältnis zueinander.

\Cref{fig:aktionen-verlauf} zeigt beide Aktionsarten im Verlauf. Die Bewegung
geht in allen vier Topologien in dieselbe Richtung, aber unterschiedlich weit.

\begin{figure}[H]
  \centering
  \includegraphics[width=\textwidth]{aktionen_verlauf}
  \caption[Reimages und Sperren im Trainingsverlauf]{Reimages und Sperren je
  Episode, gemittelt über die sechs Angreifer und fünf Seeds einer Topologie und
  geglättet über fünf Episoden. Die Kurven enden beim kürzesten der 30 Läufe,
  weil sich die Zusammensetzung des Mittels danach ändern würde.}
  \label{fig:aktionen-verlauf}
\end{figure}

Die Sperren gehen überall zurück, von mehreren hundert je Episode zu Beginn
auf 11 bis 45 am Ende. Der Rückgang ist aber ungleich: \emph{hub\_and\_spoke} und
\emph{micro\_segmented} landen bei 11 beziehungsweise 12, \emph{dmz} bei 37
und \emph{flat} bei 45. Umgekehrt steigen die
Reimages von 33 bis 51 auf 72 bis 75, mit \emph{dmz} als Ausnahme, dessen Zahl mit
36 bei etwa der Hälfte der übrigen liegt.

Der Verteidiger verlagert also sein Gewicht vom Sperren auf das Neuaufsetzen,
und zwar unabhängig von der Topologie. Bemerkenswert ist, dass ausgerechnet
\emph{flat} am Ende am häufigsten sperrt. Die naheliegende Deutung, der
Verteidiger gebe dort das Sperren auf, trifft also nicht zu.

Wie gut diese Sperren gesetzt sind, zeigt der Abwehrbonus
(\Cref{fig:abwehrbonus-verlauf}). Er ist so gebaut, dass eine Sperre gegen eine
bestehende Bedrohung positiv zählt und eine Sperre ohne Bedrohung negativ
(Gleichung~\ref{eq:abwehrbonus}); sein Vorzeichen zeigt damit an, ob die Sperren eines Laufs überwiegend
gegen eine bestehende Bedrohung standen.

\begin{figure}[H]
  \centering
  \includegraphics[width=\textwidth]{abwehrbonus_verlauf}
  \caption[Abwehrbonus im Trainingsverlauf]{Abwehrbonus je Episode, gebildet
  wie in \Cref{fig:aktionen-verlauf}. Positive Werte bedeuten, dass die
  gesetzten Sperren überwiegend gegen eine bestehende Bedrohung standen.}
  \label{fig:abwehrbonus-verlauf}
\end{figure}

Alle vier Topologien beginnen deutlich im Negativen, weil ein untrainierter
Verteidiger wahllos sperrt. Der anschließende Anstieg gegen null ist für sich
genommen aber noch kein Zeichen besseren Sperrens, denn der Bonus folgt
zunächst schlicht dem Sperrvolumen aus \Cref{fig:aktionen-verlauf}: Wo kaum
noch gesperrt wird, fällt auch kaum noch Bonus an. \emph{hub\_and\_spoke}
und \emph{micro\_segmented} erreichen die Nähe der Null mit Endniveaus von $-3{,}6$
und $-0{,}3$ auf genau diesem Weg. Sie legen das Sperren fast vollständig ab,
und die wenigen verbliebenen Sperren stehen unverändert überwiegend
nicht gegen eine bestehende Bedrohung.

Anders \emph{dmz}, dessen Endniveau mit $+0{,}6$ als einziges im Positiven
liegt. Dort pendelt sich das Sperrvolumen bei rund 30 bis 40 Sperren je
Episode ein, und diese treffen zunehmend: Ab Episode~150 fällt in gut der
Hälfte der Episoden ein positiver Bonus an, gegen drei der sechs Angreifer
ist die Bilanz im Mittel sogar positiv (\Cref{app:abwehr}). Es ist zugleich die Topologie, in
der sich das Sperren auszahlt: Vergleicht man je Matchup die stärker mit
den schwächer gesperrten Episoden nach Episode~150, so steht der
Verteidiger in den stärker gesperrten im Median 13~Punkte besser da und der Angreifer 11~Punkte schlechter. Die Zonenübergänge dieser
Topologie sind Engpässe; einer Sperre dort kann der Angreifer nicht
ausweichen. Das erklärt zugleich die Ausnahmerolle beim Neuaufsetzen aus
\Cref{fig:aktionen-verlauf}: Wer den Angreifer an den Engpässen abfängt,
hat weniger neu aufzusetzen.

\emph{flat} bleibt im Endniveau bei $-14{,}2$ stehen. Mit dem Ausreißer aus
\Cref{tab:reward-ende} hat das allerdings nichts zu tun: Ausgerechnet gegen
den eigenen Angreifer sperrt der Verteidiger ab Episode~150 im Median gar
nicht mehr, sein Bonus liegt bei null (\Cref{app:abwehr}). Der Ausreißer ist schlicht die
Grundkopplung: Der eigene Angreifer erwirtschaftet am Ende 664 Punkte
(\Cref{tab:reward-atk-verlauf}), und der Verteidiger spiegelt sie. Die Belastung stammt aus den übrigen
fünf Matchups, in denen der Verteidiger das Sperren nicht ablegt, obwohl es
sich dort nach \Cref{eq:abwehrbonus} kaum lohnen kann: Jeder Knoten in
\emph{flat} hat neun Vorgänger, und weil der Angreifer in diesen Matchups
auf zwei bis drei Knoten zurückgedrängt ist, sind fast alle davon
unkompromittiert -- eine stehende Sperre zählt damit beinahe mit dem vollen
negativen Beitrag, und genau das zeigt der anhaltend negative Bonus. Anders als in \emph{dmz} zahlt sich das nicht aus: In den stärker
gesperrten Episoden steht der Verteidiger im Median 16~Punkte schlechter,
während der Angreifer kaum gebremst wird -- die Zahl seiner gehaltenen
Knoten ändert sich nicht, denn im vollvermaschten Netz weicht er einer
Sperre schlicht auf einen anderen Knoten aus. Der Hebel des
Ratenverhältnisses aus \Cref{sec:abwehrbonus} greift hier nicht, denn er
setzt voraus, dass die Sperre einen Knoten schützt, den der Angreifer
andernfalls hielte. Warum die Policy das Verhalten dennoch beibehält, lässt
sich aus den Daten nicht bestimmen; der negative Abwehrbonus von rund 18~Punkten je Episode ist
klein gegen die Streuung des Episoden-Rewards, und entsprechend schwach ist
das Signal, das er dagegen setzt. Beide Vergleiche sind zudem Beobachtungen innerhalb der gelernten
Policies, keine kontrollierten Eingriffe.

\section{Wirksamkeit der trainierten Verteidiger}
\label{sec:wirksamkeitergebnis}

In der Evaluation (\Cref{sec:evaluation}) gibt es keinen Verlauf mehr: Je Zelle wird
erst der Median über die 25 Episoden eines Seeds gebildet, dann der Median
über die fünf Seeds; Quoten fassen alle 125 Episoden einer Zelle zusammen.
Die Stufe ohne Verteidiger liefert den Bezugspunkt für den Restanteil, der
zufällige Verteidiger trennt den Beitrag des Trainings vom zufälligen Handeln.

Der Restanteil liegt für das trainierte
Modell zwischen 0,4 und 12,1~Prozent und damit in \emph{jeder} der 24 Zellen
unter dem des Zufallsverteidigers, der zwischen 5 und 37~Prozent erreicht
(\Cref{tab:restanteil}). Das Training bewirkt also
durchgängig mehr als zufälliges Handeln, und der Abstand ist deutlich: Im
Referenzlauf hatte der zufällige Verteidiger noch nahezu gleichauf mit dem
trainierten Modell gelegen (\Cref{sec:marlon-fehler}); nach dem Umbau
trennt beide in jeder Zelle ein klarer Abstand, am größten in
\emph{micro\_segmented}: Dort liegt der Restanteil des trainierten
Modells im Mittel 16,4~Prozentpunkte unter dem des
Zufallsverteidigers. Einschränkungen
dieser Schlussfolgerung behandelt \Cref{sec:limitationen}.

\begin{table}[H]
  \centering
  \caption[Restanteil in der Evaluation]{Restanteil des Angreifer-Rewards
  gegen den trainierten Verteidiger, in
  Prozent des ungeschützten Netzes; dieses selbst liegt definitionsgemäß bei
  100. In Klammern der Wert des Zufallsverteidigers. Niedriger ist besser.}
  \label{tab:restanteil}
  \small
  \begin{tabular}{l*{6}{r}}
    \toprule
    Verteidigt & Chain & Flat & H\,\&\,S & DMZ & Micro & Super \\
    \midrule
    Flat         & 8,8 (14) & 12,1 (19) & 9,0 (16) & 7,6 (12) & 8,7 (15) & 3,0 (9)  \\
    Hub \& Spoke & 5,5 (11) & 6,4 (15)  & 0,5 (10) & 3,9 (6)  & 9,8 (19) & 0,4 (5)  \\
    DMZ          & 4,1 (9)  & 3,1 (7)   & 3,8 (10) & 2,4 (10) & 4,7 (9)  & 3,6 (7)  \\
    Micro-Seg.   & 4,3 (11) & 4,9 (16)  & 3,0 (22) & 10,2 (12)& 1,2 (37) & 3,6 (27) \\
    \bottomrule
  \end{tabular}
\end{table}

Die Crown-Jewel-Quote bestätigt das Bild und schärft es
(\Cref{tab:cj-eval}). Gegen das trainierte Modell erreicht der Angreifer das Crown Jewel in
nahezu allen Zellen deutlich seltener als gegen den Zufallsverteidiger. In
\emph{micro\_segmented} fällt die Quote bei drei von sechs Angreifern auf
exakt null -- beim eigenen Angreifer, der gegen den Zufallsverteidiger noch
jede Episode gewann; der Chain- und der Flat-Angreifer kommen in je einer
einzigen der 125 Episoden einer Zelle durch.
Das andere Extrem ist erneut \emph{flat} gegen den eigenen Angreifer mit
99,2~Prozent: Hier erreicht der Angreifer das Crown Jewel praktisch in jeder
Episode, obwohl ein trainierter Verteidiger im Netz steht.

\begin{table}[H]
  \centering
  \caption[Crown-Jewel-Quote in der Evaluation]{Anteil der Episoden mit
  erreichtem Datenbankserver gegen den trainierten Verteidiger, in Prozent.
  Die kleine Zeile darunter nennt Zufallsverteidiger und ungeschütztes
  Netz.}
  \label{tab:cj-eval}
  \footnotesize
  \setlength{\tabcolsep}{4pt}
  \begin{adjustbox}{max width=\textwidth}
  \begin{tabular}{l*{6}{r}}
    \toprule
    Verteidigt & Chain & Flat & H\,\&\,S & DMZ & Micro & Super \\
    \midrule
    Flat         & 51,2 & 99,2 & 53,6 & 25,6 & 49,6 & 84,0 \\
    \multicolumn{1}{r}{\scriptsize Zufall\,/\,ohne} & \scriptsize 92\,/\,92 & \scriptsize 100\,/\,100 & \scriptsize 94\,/\,100 & \scriptsize 55\,/\,66 & \scriptsize 90\,/\,95 & \scriptsize 100\,/\,100 \\[2pt]
    Hub \& Spoke & 8,8 & 4,0 & 27,2 & 2,4 & 20,0 & 2,4 \\
    \multicolumn{1}{r}{\scriptsize Zufall\,/\,ohne} & \scriptsize 46\,/\,68 & \scriptsize 50\,/\,77 & \scriptsize 100\,/\,100 & \scriptsize 14\,/\,94 & \scriptsize 54\,/\,69 & \scriptsize 57\,/\,90 \\[2pt]
    DMZ          & 8,0 & 0,8 & 1,6 & 36,0 & 0,8 & 1,6 \\
    \multicolumn{1}{r}{\scriptsize Zufall\,/\,ohne} & \scriptsize 46\,/\,78 & \scriptsize 43\,/\,98 & \scriptsize 58\,/\,95 & \scriptsize 99\,/\,100 & \scriptsize 31\,/\,98 & \scriptsize 25\,/\,99 \\[2pt]
    Micro-Seg.   & 0,8 & 0,8 & 0,0 & 0,0 & 0,0 & 15,2 \\
    \multicolumn{1}{r}{\scriptsize Zufall\,/\,ohne} & \scriptsize 39\,/\,100 & \scriptsize 52\,/\,100 & \scriptsize 6\,/\,56 & \scriptsize 0\,/\,44 & \scriptsize 100\,/\,100 & \scriptsize 100\,/\,100 \\[2pt]
    \bottomrule
  \end{tabular}
  \end{adjustbox}
\end{table}

Die Zahl der gehaltenen Knoten weist in dieselbe Richtung
(\Cref{tab:knoten-eval}). Gegen das trainierte Modell hält der Angreifer im
Median 2 von zehn Knoten, den Einstiegsknoten mitgezählt, also faktisch
einen einzigen eroberten. Ohne Verteidiger sind es je nach Zelle 3 bis 10,
gegen den Zufallsagenten 2 bis 8. Ausnahmen beim trainierten Modell sind wiederum \emph{flat} gegen \emph{flat} mit 7
Knoten sowie \emph{flat} gegen den Super-Angreifer und \emph{dmz} gegen \emph{dmz} mit je 3.

\begin{table}[H]
  \centering
  \caption[Gehaltene Knoten in der Evaluation]{Größte Zahl gleichzeitig
  gehaltener Knoten gegen den trainierten Verteidiger, als Median einer
  Zelle; die kleine Zeile darunter nennt Zufallsverteidiger und
  ungeschütztes Netz. Von zehn Knoten, der Einstiegsknoten zählt mit.}
  \label{tab:knoten-eval}
  \footnotesize
  \setlength{\tabcolsep}{4pt}
  \begin{adjustbox}{max width=\textwidth}
  \begin{tabular}{l*{6}{r}}
    \toprule
    Verteidigt & Chain & Flat & H\,\&\,S & DMZ & Micro & Super \\
    \midrule
    Flat         & 2 & 7 & 2 & 2 & 2 & 3 \\
    \multicolumn{1}{r}{\scriptsize Zufall\,/\,ohne} & \scriptsize 3\,/\,7 & \scriptsize 7\,/\,10 & \scriptsize 3\,/\,7 & \scriptsize 3\,/\,9 & \scriptsize 3\,/\,8 & \scriptsize 4\,/\,10 \\[2pt]
    Hub \& Spoke & 2 & 2 & 2 & 2 & 2 & 2 \\
    \multicolumn{1}{r}{\scriptsize Zufall\,/\,ohne} & \scriptsize 3\,/\,8 & \scriptsize 3\,/\,8 & \scriptsize 8\,/\,10 & \scriptsize 3\,/\,9 & \scriptsize 3\,/\,3 & \scriptsize 3\,/\,10 \\[2pt]
    DMZ          & 2 & 2 & 2 & 3 & 2 & 2 \\
    \multicolumn{1}{r}{\scriptsize Zufall\,/\,ohne} & \scriptsize 3\,/\,8 & \scriptsize 3\,/\,8 & \scriptsize 3\,/\,8 & \scriptsize 8\,/\,10 & \scriptsize 3\,/\,5 & \scriptsize 3\,/\,9 \\[2pt]
    Micro-Seg.   & 2 & 2 & 2 & 2 & 2 & 2 \\
    \multicolumn{1}{r}{\scriptsize Zufall\,/\,ohne} & \scriptsize 3\,/\,4 & \scriptsize 3\,/\,4 & \scriptsize 3\,/\,4 & \scriptsize 2\,/\,3 & \scriptsize 4\,/\,4 & \scriptsize 4\,/\,4 \\[2pt]
    \bottomrule
  \end{tabular}
  \end{adjustbox}
\end{table}

Beide Größen zusammen zeigen, dass die drei segmentierten Topologien einander
im Ergebnis ähnlicher sind, als der Reward vermuten ließ, und dass \emph{flat}
klar abfällt. Bemerkenswert ist dabei, dass Restanteil und Crown-Jewel-Quote
nicht dieselbe Rangfolge liefern. \emph{hub\_and\_spoke} gegen den eigenen Angreifer hat
mit 0,5~Prozent einen der niedrigsten Restanteile der Tabelle, zugleich
aber eine Crown-Jewel-Quote von 27,2~Prozent. Der Restanteil ist ein Verhältnis zum
ungeschützten Netz; wo dort sehr viel zu holen war, sieht eine Verteidigung relativ
gut aus, obwohl das Crown Jewel in mehr als jeder vierten Episode fällt. Für
eine Aussage über das tatsächliche Risiko ist deshalb die absolute Größe
maßgeblich.

Schließlich die Gegenrechnung. \Cref{tab:kosten} stellt den Aufwand beider
Verteidiger gegenüber. Der Zufallsagent setzt im Median rund 650 Sperren
je Episode, also bei etwa jedem dritten der 2000 Schritte, und verletzt
dadurch in 99,7 bis 100~Prozent der Episoden das \ac{SLA}; sein Reward liegt
entsprechend um das Vierzig- bis Hundertsechzigfache unter dem des
trainierten Modells. Dieses sperrt fast nie und bricht das \ac{SLA} nur
selten. Der Unterschied liegt damit nicht in der Menge der Eingriffe, sondern
in ihrer Auswahl.

\begin{table}[H]
  \centering
  \caption[Aufwand der Verteidigung]{Aufwand des trainierten Verteidigers je
  Evaluationsepisode, in Klammern der Zufallsverteidiger. Mediane über alle
  Seeds und Angreifer.}
  \label{tab:kosten}
  \small
  \begin{tabular}{lrrrr}
    \toprule
    Verteidigt & Sperren & Reimages & \ac{SLA}-Bruch & Reward \\
    \midrule
    Flat         & 10 (650) & 29,0 (21,0)  & 14,9\,\% (100\,\%)  & $-191$ ($-8016$) \\
    Hub \& Spoke & 1 (650)  & 103,0 (34,0) & 0,8\,\% (100\,\%)   & $-91$ ($-8356$)  \\
    DMZ          & 19 (654) & 27,0 (16,0)  & 4,5\,\% (99,7\,\%)  & $-78$ ($-6934$)  \\
    Micro-Seg.   & 0 (648)  & 103,5 (34,0) & 0,1\,\% (100\,\%)   & $-53$ ($-8497$)  \\
    \bottomrule
  \end{tabular}
\end{table}

Auffällig ist die Zeile \emph{micro\_segmented}: null Sperren im Median, dafür 103,5
Reimages und damit dreimal so viele wie beim Zufallsagenten. Der wirksamste
Verteidiger dieser Untersuchung sperrt also gar nicht, sondern setzt konsequent
neu auf. Dasselbe Muster zeigt \emph{hub\_and\_spoke}. \Cref{sec:rewarddiskussion} geht
der Frage nach, ob das eine Eigenschaft der Topologien ist oder eine des
Reward-Entwurfs.

\section{Einfluss der Angreiferherkunft}
\label{sec:angreiferherkunft}

Die Topologie, auf der ein Angreifer trainiert wurde, beeinflusst das Ergebnis
sichtbar. In drei der vier verteidigten Topologien erreicht derjenige Angreifer
die höchste Crown-Jewel-Quote, der auf ebendieser Topologie entstanden ist:
\emph{flat} mit 99,2~Prozent, \emph{hub\_and\_spoke} mit 27,2 und \emph{dmz} mit 36,0. Wie deutlich
der Abstand zur zweithöchsten Quote derselben Zeile ausfällt, ist
allerdings sehr verschieden (\Cref{tab:cj-eval}): In \emph{dmz} folgt der Chain-Angreifer erst bei 8,0~Prozent, in
\emph{hub\_and\_spoke} der Micro-Segmentation-Angreifer bei 20,0, und in \emph{flat} liegt der
Super-Angreifer bei 84,0.

Der Vorteil ist also kein allgemeiner Fähigkeitsvorsprung, sondern
Passgenauigkeit: Der Angreifer kennt die Wege dieser Struktur, weil er auf ihr
gelernt hat. Umgekehrt bricht seine Leistung ein, sobald er auf eine fremde und
stärker segmentierte Struktur trifft.

Die Ausnahme ist \emph{micro\_segmented} selbst. Dort scheitert der eigene
Angreifer vollständig (0,0~Prozent), während der Super-Angreifer als einziger
mit 15,2~Prozent durchkommt. Die Vielfalt seiner
Trainingsstationen vor der letzten auf \emph{micro\_segmented} verschafft
ihm hier offenbar einen Vorteil gegenüber dem Spezialisten: Wo eine
Struktur nur wenige und enge Wege lässt, scheint breites Wissen mehr zu
helfen als spezialisiertes.

In einzelnen Konstellationen kann breites Wissen für einen Verteidiger
also gefährlicher sein als der Spezialist. Für drei der vier Topologien
gilt jedoch das Gegenteil: Dort liefert der eigens für die verteidigte
Topologie trainierte Angreifer die besseren Ergebnisse, und der Generalist
bleibt teils weit zurück -- gegen die Verteidiger von \emph{hub\_and\_spoke} und \emph{dmz}
etwa mit 2,4 und 1,6~Prozent gegenüber 27,2 und 36,0~Prozent der
Spezialisten (\Cref{tab:cj-eval}). Auch die Ausnahme trägt zudem einen
Vorbehalt: Das gestufte Training des Super-Angreifers endete auf
\emph{micro\_segmented} (\Cref{sec:phasen}); seine Stärke gerade dort kann
deshalb ebenso Aktualität wie Breite sein -- die Policy wurde zuletzt auf
dieser Topologie angepasst, und früher Gelerntes wird dabei teilweise
überschrieben. Die vorliegenden Daten lassen beide Deutungen zu; keine von ihnen
lässt sich ausschließen.

Für die Fragestellung folgt daraus zweierlei. Erstens ist die
Angreiferherkunft eine Einflussgröße eigener Art und darf nicht über alle
Angreifer hinweg gemittelt werden, wie in \Cref{sec:evaluation} begründet.
Zweitens beschreibt die Diagonale den ungünstigsten Fall, den ein Verteidiger
zu erwarten hat, nämlich einen Gegner, der genau diese Topologie kennt.

\section{Einordnung des Reward-Designs}
\label{sec:rewarddiskussion}

Der auffälligste Befund dieses Kapitels ist, dass die wirksamsten Verteidiger
kaum sperren. Diese Beobachtung lässt zwei Deutungen zu; anhand der erhobenen
Daten lässt sich zwischen ihnen nicht abschließend entscheiden.

Die erste Deutung ist inhaltlich: In einem Netz, in dem ohnehin nur wenige
Wege bestehen, ist das Neuaufsetzen das schärfere Mittel, weil es den
Angreifer entfernt, statt ihn nur aufzuhalten. Dafür spricht, dass gerade die
Verteidiger der beiden Topologien mit den wenigsten Kanten am häufigsten neu
aufsetzen.

Die zweite Deutung liegt im Reward-Entwurf. Der Abwehrbonus ist mit
$0{,}00025 \cdot \text{Wert}(Z)$ je Schritt bewusst klein gehalten
(\Cref{sec:abwehrbonus}), damit das Neuaufsetzen stets lohnender bleibt; das
Verhältnis der Raten beträgt 20. Ein Reimage bringt dagegen sofort den vollen
Knotenwert. Prävention ist damit zwar sichtbar, aber deutlich schwächer
vergütet als Wiederherstellung, und der Agent folgt dieser Gewichtung.

Beide Deutungen sind mit den Daten vereinbar. Zu trennen wären sie nur durch
Läufe mit systematisch veränderter Abwehrbonusrate, die in dieser Arbeit nicht
durchgeführt wurden. Festzuhalten bleibt, dass die in
\Cref{sec:mehrstufig} beschriebene Einschließungsstrategie in keinem Lauf
auftrat: Der Abwehrbonus genügte nicht, um eine mehrstufige Sperrkette
aufzubauen.

Ob sich damit auch die Diagnose des Referenzlaufs bestätigt, nach der
Prävention im Reward unsichtbar blieb (\Cref{sec:marlon-fehler}), lässt
sich dagegen nicht sagen: Die beiden Läufe sind dafür zu verschieden
(\Cref{sec:limitationen}). Eine Kontrolle des Entwurfs liefert dafür die Obergrenze aus der
Grundkopplung: Der Abwehrbonus ist so klein bemessen, dass er sich nicht zu
positivem Reward farmen lassen soll (\Cref{sec:abwehrbonus}) -- und
tatsächlich bleibt der Verteidiger-Reward in jeder einzelnen Episode von
Training und Evaluation negativ, der beste gemessene Einzelwert liegt bei
$-12{,}8$.

\section{Limitationen}
\label{sec:limitationen}

Das Konvergenzkriterium unterscheidet die Topologien nicht. Es ist eigens für
diese Arbeit entworfen und nicht durch Literatur gedeckt
(\Cref{sec:konvergenz}); seine Schwäche liegt in der Normierung auf die
Gesamtspanne des Laufs, die von den extremen Startwerten der Anfangsphase
bestimmt wird. Ein strengerer Trendtest verschiebt den Zeitpunkt, ändert an
dieser Bezugsgröße aber nichts (\Cref{sec:konvergenzergebnis}).

Ein Maß ohne diese Schwäche ließe sich so bilden: Man bestimmt je Lauf den
Median seiner ersten und seiner letzten zehn Episoden und bestimmt, ab
welcher
Episode er 90~Prozent der Strecke dazwischen erreicht hat und zehn Episoden
lang hält. Bezugspunkt ist damit die Verbesserung des Laufs selbst, nicht seine
Gesamtspanne. Gesucht wird erst ab Episode~11, weil kein Lauf als konvergiert
gelten kann, bevor sein eigener Bezugspunkt feststeht; der frühestmögliche Wert
liegt damit bei Episode~20.

Angewendet auf die vorliegenden Daten erreichen alle 120 Läufe diese Marke, im
Median in Episode~43. Die Mediane je Topologie lauten 50 für \emph{flat}, 45
für \emph{dmz}, 42 für \emph{hub\_and\_spoke} und 39 für
\emph{micro\_segmented} und ordnen die Topologien damit genauso wie Reward,
Crown-Jewel-Quote und gehaltene Knoten. Nachweisbar ist auch diese Ordnung
nicht, die Quartile reichen von 38 bis 68 in \emph{flat} und von 31 bis 55 in
\emph{micro\_segmented}. Der Befund aus \Cref{sec:konvergenzergebnis} hängt also
weder an der Schwelle noch an der Konstruktion des Maßes, sondern ebenso an der
Zahl der Läufe. Weitere naheliegende Schritte wären ein Ausschluss der
Anfangsphase und vor allem mehr Seeds je Zelle; eine Prüfung dieser
Varianten liegt jedoch außerhalb des Umfangs dieser Arbeit.

Jede Zelle beruht auf fünf Seeds. Henderson et~al.\ zeigen, dass allein der
Zufallsstartwert bei Policy-Gradient-Verfahren zu deutlich verschiedenen
Ergebnisverteilungen führen kann, und raten deshalb zu mehr Läufen und zur
Angabe von Konfidenzbereichen~\cite{henderson2019deepreinforcementlearningmatters}.
Fünf Seeds sind das untere Ende des dort empfohlenen Wertebereichs. Abgemildert wird das an
zwei Stellen: Es werden alle Läufe berichtet und keine ausgewählt, und die
Evaluation stützt sich auf Mediane, die ein einzelner Ausreißer nicht
verschiebt. Für die Lernkurven gilt das nicht, dort schlägt ein einzelner Seed
in das Mittel durch, wie \Cref{sec:lernkurven} zeigt. Die berichteten
Unterschiede zwischen den Topologien sind allerdings so groß, dass sie sich
nicht durch Streuung zwischen Seeds erklären lassen; kleinere Unterschiede, etwa
zwischen \emph{dmz} und \emph{micro\_segmented}, wären erst mit einer
größeren Zahl von Läufen aussagekräftig.

Der Angreifer ist während des Verteidigertrainings eingefroren
(\Cref{sec:phasen}). Der Verteidiger kann sich deshalb auf ein festes
Verhalten einspielen, und die gemessene Wirksamkeit gilt zunächst nur gegen
diesen einen Gegner. Die Auswertung je Matchup mildert diesen Umstand, da sechs
verschiedene Angreifer geprüft werden, hebt ihn aber nicht
vollständig auf. Das
relativiert insbesondere den Befund, dass der passgenaue Angreifer am
Trainingsende neutralisiert ist: Er gelingt einem Verteidiger, der sich
ungestört auf genau dieses eine Verhalten einspielen konnte. Gegen einen
mitlernenden Angreifer stünde diese Möglichkeit nicht offen, und der
Einfluss der Angreiferherkunft könnte entsprechend größer ausfallen.

Auch die Reihenfolge, in der der Super-Angreifer die fünf Topologien
durchlief, wurde nicht variiert. Sie endete auf \emph{micro\_segmented}, und
\Cref{sec:angreiferherkunft} zeigt, dass sich seine Stärke dort nicht von
einem bloßen Aktualitätsvorteil der letzten Station trennen lässt.

Der Hub in \emph{hub\_and\_spoke} ist ein übernehmbarer Knoten
(\Cref{sec:topomodell}). In einem realen Netz wäre die Vermittlungsstelle
häufig ein Gerät ohne Nutzerdienste. Die Ergebnisse dieser Topologie sind
deshalb nicht ohne Weiteres auf eine Sterntopologie mit gehärtetem Zentrum zu
übertragen.

Zwei Entwurfsentscheidungen am Aktionsraum wurden nicht validiert. Der
Dienststopp wurde entfernt, weil die knotenweite Sperre dieselbe Wirkung
erzielt; ein Lauf mit beiden Aktionen nebeneinander wurde nicht gespielt
(\Cref{sec:aktionsraum-def}). Ebenso wurde die Sperre
von portweise auf knotenweise umgestellt, ohne beide Varianten zu
vergleichen; an dieser Entscheidung hängt auch die Deutung der Sperrwirkung
in \Cref{sec:aktionen}, denn eine portweise Sperre träfe in \emph{flat}
nicht alle neun Zugänge zugleich. Ungeprüft blieb schließlich, ob der Verteidiger in \emph{flat}
ohne sein dauerhaftes, überwiegend negativ verbuchtes Sperren besser
abschnitte (\Cref{sec:aktionen}); auch dies würde eigene Läufe erfordern.

Eine Lücke betrifft schließlich die Erhebung selbst. Die Aktionsstatistik
zählt Freigaben allein nach Aktionstyp und hält nicht fest, ob der betroffene
Knoten zum Zeitpunkt der Freigabe überhaupt gesperrt war. Berechtigte
Freigaben lassen sich deshalb nachträglich nicht von folgenlosen trennen, was
die Deutung ihres Anteils in \Cref{sec:aktionen} begrenzt. Der Sperrzustand
liegt zur Laufzeit vor; ihn mitzuschreiben wäre der naheliegende Weg, diese
Lücke zu schließen.

Aus dem Vergleich mit dem Referenzlauf lassen sich nur eingeschränkt
Ursachen ableiten.
Zwischen beiden Läufen wurden mehrere Änderungen gleichzeitig vorgenommen,
nämlich der Abwehrbonus, der Aktionsraum, die Maskierung, drei
Fehlerkorrekturen -- und der Wegfall des Abbruchkriteriums, das die
Referenzläufe schon nach wenigen Dutzend Episoden beendete
(\Cref{sec:referenzlauf}). Welcher Anteil der Verbesserung auf welche
Änderung entfällt, lässt sich daraus nicht bestimmen; insbesondere kann die
Nähe von trainiertem und zufälligem Verteidiger im Referenzlauf auch an
dessen kurzem Training gelegen haben, nicht allein an der im Reward
unsichtbaren Prävention.

Schließlich bleibt der Aufbau eine Simulation. \ac{CBS} bildet Angriffe als
Folge diskreter Schritte auf einem gerichteten Graphen ab, ohne Zeitverhalten,
ohne Erkennungsunsicherheit und mit einem Verteidiger, der den Zustand seines
Netzes vollständig kennt. Die Ergebnisse sind Aussagen über dieses Modell. Ob
sich die beobachtete Rangfolge der Topologien auf reale Netze überträgt, ist
damit nicht gezeigt.
